\setchapter{4}{Limited Dependent Variable Models}
\ChapterHeading

In some applications, the domain of the variate of interest is limited.
In application, the bounds of the domain of $Y$ often have to be taken into account.
Traditionally, there are three forms of limited dependent variable models:
\begin{enumerate}
    \item Truncated data
    \item Censored data
    \item Sample selection which is also called incidental truncation
\end{enumerate}

\section{Truncated data}

\begin{definition}
    A sample $Y$ is said to come from a \textbf{truncated distribution} when it is not possible to obtain observations from part of the domain of $Y$.
\end{definition}

If $Y$ is a \Emph{continuous random variable} with pdf $f(y)$ and $a$ is a constant, we have \Emph{left-truncation at $a$} if the sample is drawn from the conditional distribution 
\begin{equation*}
    f(y|Y>a) = \frac{f(y)}{P(Y>a)}
\end{equation*}
We have also truncation if the regressors are also not observed when $Y<a$.
For a \Emph{discrete random variable} with support on $0,1,\ldots,\infty$, we have \Emph{left-truncation at $a$} if the sample is drawn from 
\begin{equation*}
    P(Y=j|Y>a) = \frac{P(Y=j)}{1-\sum_{k=0}^{\lfloor a \rfloor}P(Y=k)'}
\end{equation*}

Notice that truncation is only problematic if it depends on $Y$ or on another variable that is not independent of $Y$ and that is not used as a regressor.
Under certain conditions, it is possible to perform inference about the entire population using truncated samples.
The problem is that this inference tends to be \Emph{sensitive to distributional assumptions}.

\begin{Example}
    Normal distribution with mean $\mu$ and variance $\sigma^2$
    \begin{equation*}
        E[Y_i|Y_i>a] = \mu + \underbrace{\sigma\frac{\phi(\frac{1-\mu}{\sigma})}{1-\Phi(\frac{a-\mu}{\sigma})}}_{\text{extra term}}
    \end{equation*}
    If you try to make inference on $\mu$, it would be biased because of this extra term.\\

    For the case of count data with $E[Y_i|X_i] = \exp(X'_i\beta_0)$
    \begin{equation*}
        E[Y_i|X_i,Y_i>a] = \frac{\exp(X'_i\beta_0)}{1 - \underbrace{\sum_{k=0}^{\lfloor a \rfloor}P(Y_i=k|X_i)}_{\text{notice here}}}
    \end{equation*}
    The result depends on the probability function of the data.
\end{Example}
In both of the examples, the expectation for the truncated data \Emph{depends on the shape of the conditional distribution}.
Because of this, the estimation is often performed by \KeyIdea{maximum likelihood}.
\Info{Assume we know the distribution of the data.}

\subsection{Truncated normal regression model}
The \Term{Truncated normal regression model} is defined by 
\begin{align*}
    Y_i^* & = X'_i\beta_0+u_i^*,\quad u_i^*\sim\mcN(0,\sigma^2)\\
    (Y_i,X_i) & = \begin{cases}
        (Y_i^*,X_i) & \text{if }Y_i^* > 0\\
        \text{non-observable} & \text{if }Y_i^*\leq 0
    \end{cases}
\end{align*}
This implies that 
\begin{equation*}
    E[Y_i|X_i, Y_i>0] = X'_i\beta_0 + \sigma\frac{\phi(\frac{-X'_i\beta_0}{\sigma})}{1-\Phi(\frac{-X'_i\beta_0}{\sigma})}
\end{equation*}
\Info{Again with the extra term, the OLS assumption of $E[Y_i|X_i]=X_i'\beta_0$ is violated $\to$ use MLE.}
The parameters of interest can be estimated from the likelihood function 
\begin{equation*}
    L(\beta,\sigma) = \prod_{i=1}^{n}\frac{1}{\sigma}\frac{\phi(\frac{Y_i-X'_i\beta}{\sigma})}{1-\Phi(\frac{-X'_i\beta_0}{\sigma})}
\end{equation*}

\section{Censored data}

Censoring is a problem of partial observability.
For example, if $Y^*$ is a random variable and $a$ is a constant, we have \Emph{left-censoring at $a$} if we can only observe 
\begin{equation*}
    Y = \max\{Y^*,a\} = \begin{cases}
        a & \text{if }Y^* \leq a\\
        Y^* & \text{if }Y^* > a
    \end{cases}
\end{equation*}
We can also have \Emph{right truncation}, in this case, $Y = \min\{Y^*,a\}$.\\

Focusing on left-censored data.
If $Y^*$ is a \Emph{discrete random variable} with support on $0,1,\ldots,\infty$, the probability function of $Y$ is given by 
\begin{equation*}
    P(Y=j) = \begin{cases}
        0 & \text{if }j<a\\
        P(Y^*\leq a) & \text{if }j=a\\
        P(Y^* = j) & \text{if }j>a
    \end{cases}
\end{equation*}
If $Y^*$ is a \Emph{continuous random variable} with pdf $f(y^*)$, then $Y$ has the \Emph{mixed distribution} 
\begin{equation*}
    P(Y\leq k) = \begin{cases}
        0 & \text{if }k<a\\
        P(Y^*\leq k) & \text{if }k\geq a
    \end{cases}
\end{equation*}
It is assumed that the \Emph{regressors are fully observed}, even for cases with censored $Y^*$.

\subsection{Tobit model}

The standard regression model for continuous censored data is the \Term{Tobit} which was proposed by Tobin (1958).
The model is defined by 
\begin{align*}
    Y^*_i & = X'_i\beta_0 + u_i^*,\quad u_i^*\sim\mcN(0,\sigma^2)\\
    Y_i & = \begin{cases}
        0 & \text{if }Y_i^* \leq 0\\
        Y_i & \text{if }Y_i^*>0
    \end{cases}
\end{align*}
The model is appropriate if we are interested in \Emph{making inference about $Y_i^*$ not $Y_i$}.
In this case,
\begin{equation*}
    E[Y_i|X_i] = X'_i\beta_0\Phi(\frac{X'_i\beta_0}{\sigma})+\sigma\phi(\frac{X'_i\beta_0}{\sigma})
\end{equation*}
\KeyIdea{ML inference} is based on the likelihood function 
\begin{equation*}
    L(\beta,\sigma) = \prod_{i=1}^{n}\left[1-\Phi(\frac{X'_i\beta_0}{\sigma})\right]^{1(Y_i=0)}\left[\frac{1}{\sigma}\phi\left(\frac{y_i-X'_i\beta}{\sigma}\right)\right]^{1(Y_i>0)}
\end{equation*}



% \begin{definition}
% Theorem
% \begin{lemma}
% \begin{Example}
% \begin{proof}
%\newcommand{\KeyIdea}  % 关键想法：浅绿底
%\newcommand{\Term}     % 专业名词：品牌深绿字
%\newcommand{\Emph}     % 强调词汇：红
%\newcommand{\Confuse}    % 不懂/存疑处
%\newcommand{\Info}     % 说明或注解